% Testing Document — Vehicle Dynamics (IEEEtran)
% Requires IEEEtran.cls (TeX Live / MiKTeX IEEE package).

\documentclass[conference]{IEEEtran}
\usepackage{amsmath,amssymb,amsfonts}
\usepackage{booktabs}
\usepackage{cite}
\usepackage{url}
\usepackage{hyperref}

\begin{document}

\title{Software Testing Document\\for a Vehicle Dynamics Simulator with DEM Soft Soil}

\author{\IEEEauthorblockN{Sankalp Amaravadi}
\IEEEauthorblockA{Vehicle Dynamics Project\\
Email: (TBD)}}

\maketitle

\begin{abstract}
This testing document defines strategy, levels, and initial cases for the
Vehicle Dynamics backend. Coverage spans unit, integration, and system tests for
orchestration, vehicle models, DEM soil, and tire--soil coupling. GPU-scale
performance tests are treated as optional, environment-specific checks. User
interface testing is deferred with the front end.
\end{abstract}

\begin{IEEEkeywords}
software testing, verification, vehicle dynamics, DEM, pytest, GPU
\end{IEEEkeywords}

\section{Introduction}

\subsection{Purpose}
Describe how correctness, regression safety, and (where applicable) performance
are assessed for the Vehicle Dynamics simulator.

\subsection{Scope}
In scope: backend modules and end-to-end runs via \texttt{main.py}. Out of
scope: GUI testing until a front end exists.

\subsection{References}
SRS (\texttt{srs.tex}) and SDD (\texttt{sdd.tex}) define requirements and
design under test.

\section{Test Strategy}
\begin{itemize}
  \item Unit tests for pure Python modules (configuration, helpers, I/O).
  \item Integration tests for vehicle--DEM coupling with small particle counts.
  \item Optional GPU performance / scale checks for large DEM runs (non-blocking
        for default CI).
  \item Regression scenarios with golden-output tolerances (TBD).
\end{itemize}

\section{Test Levels}

\subsection{Unit Testing}
Verify individual functions and classes in isolation with deterministic inputs.

\subsection{Integration Testing}
Verify tire--soil coupling and time-loop ordering using tiny DEM fields that
run on CPU or modest GPU allocations.

\subsection{System Testing}
Execute end-to-end scenarios through \texttt{main.py} and confirm expected
artifacts and exit status.

\subsection{Acceptance / Performance (Optional)}
Measure wall-clock time and memory for large particle counts on GPU hardware.
Results are informative unless a quantitative NFR threshold is adopted.

\section{Test Cases}
\begin{itemize}
  \item \textbf{TC-1:} \texttt{main.py} completes successfully (Phase~1 smoke).
  \item \textbf{TC-2:} Configuration loading rejects invalid parameters (TBD).
  \item \textbf{TC-3:} A small DEM step satisfies stated invariants within
        tolerance (TBD).
  \item \textbf{TC-4:} Exported time series contain required columns and
        monotonically increasing time (TBD).
\end{itemize}

\section{Environment and Tooling}
\begin{itemize}
  \item Language runtime: Python~3.
  \item Planned automation: \texttt{pytest}.
  \item GPU tests: compatible NVIDIA drivers and selected DEM/GPU stack when
        scale tests are enabled.
\end{itemize}

\section{Pass / Fail Criteria}
A build passes default CI when all mandatory unit, integration, and system
smoke tests succeed. Optional GPU-scale jobs may fail independently without
blocking documentation-only or CPU-smoke pipelines until NFR thresholds are
frozen.

\section*{Acknowledgment}
This document follows the IEEE conference \LaTeX\ template (\texttt{IEEEtran}).

\begin{thebibliography}{00}
\bibitem{ieee829} IEEE Std 829-2008, \emph{IEEE Standard for Software and
System Test Documentation}.
\bibitem{ieeetran} M.~Shell, ``IEEEtran homepage,''
\url{https://www.michaelshell.org/tex/ieeetran/}.
\end{thebibliography}

\end{document}
